\textbf{R:} Comencemos por recordar el modelo de regresion lineal
\[ y_{pred} = \sum_{i=1}^{d} w_i x_i + b\]
Y las reglas de actualización:
\begin{align*}
    error &= y-y_{pred} \\
    w_j &\gets w_j + \eta \cdot error \cdot x_j \text{ Para cada } j\\
    b &\gets b + \eta \cdot error
\end{align*}

\subsection*{Iteración 1}
Comenzamos con el estado inicial indicado que es de $w_1 = w_2 = 1$ y $b=1$.
\subsubsection*{Muestra 1}
Tenemos $x_1 = 0$, $x_2 = 0$ y $y=0$, tenemos
\begin{align*}
    y_{pred} &= 1\cdot 0 + 1 \cdot 0 + 1 = 1 \\
    error &= 0 - 1 = -1 \\
    w_1 &= 1 + 0.5 \cdot (-1) \cdot 0 = 1 \\
    w_2 &= 1 + 0.5 \cdot (-1) \cdot 0 = 1 \\
    b &= 1 + 0.5 \cdot (-1) = 0.5
\end{align*}
Se tienen los valores actualizados $w(1 \quad 1), b= 0.5$.
    
\subsubsection*{Muestra 2}
Tenemos $x_1 = 0$, $x_2 = 1$ y $y=0$, así que 
\begin{align*}
    y_{pred} &= 1\cdot 0 + 1 \cdot 1 + 0.5 = 1.5 \\
    error &= 0 - 1.5 = -1.5 \\
    w_1 &= 1 + 0.5 \cdot (-1.5) \cdot 0 = 1 \\
    w_2 &= 1 + 0.5 \cdot (-1.5) \cdot 1 = 0.25 \\
    b &= 0.5 + 0.5 \cdot (-1.5) = -0.25
\end{align*}
Se tienen los valores actualizados $w(1 \quad 0.25), b= -0.25$.

\subsubsection*{Muestra 3}
Tenemos $x_1 = 1$, $x_2 = 0$ y $y=0$, tenemos
\begin{align*}
    y_{pred} &= 1 \cdot 1 + 0.25 \cdot 0 - 0.25 = 0.75 \\
    error &= 0 - 0.75 = -0.75 \\
    w_1 &= 1 + 0.5 \cdot (-0.75) \cdot 1 = 0.625 \\
    w_2 &= 0.25 + 0.5 \cdot (-0.75) \cdot 0 = 0.25 \\
    b &= -0.25 + 0.5 \cdot (-0.75) = -0.625
\end{align*}

Se tienen los valores actualizados $w(0.625 \quad 0.25), b= -0.625$.

\subsubsection*{Muestra 4}
Tenemos $x_1 = 1$, $x_2 = 1$ y $y=1$, tenemos
\begin{align*}
    y_{pred} &= 0.625 \cdot 1 + 0.25 \cdot 1 - 0.625 = 0.25 \\
    error &= 1 - 0.25 = 0.75 \\
    w_1 &= 0.625 + 0.5 \cdot (0.75) \cdot 1 = 1 \\
    w_2 &= 0.25 + 0.5 \cdot (0.75) \cdot 1 = 0.625 \\
    b &= -0.625 + 0.5 \cdot (0.75) = -0.25
\end{align*}
La iteración 1 finaliza con los valores $w(1 \quad 0.625), b= -0.25$.


\subsection*{Iteración 2}
\subsubsection*{Muestra 1}
Tenemos $x_1 = 0$, $x_2 = 0$ y $y=0$, tenemos
\begin{align*}
    y_{pred} &= 1\cdot 0 + 0.625 \cdot 0 - 0.25 = -0.25 \\
    error &= 0 - (-0.25) = 0.25 \\
    w_1 &= 1 + 0.5 \cdot (0.25) \cdot 0 = 1 \\
    w_2 &= 0.625 + 0.5 \cdot (0.25) \cdot 0 = 0.625 \\
    b &= -0.25 + 0.5 \cdot (0.25) = -0.125
\end{align*}
Se tienen los valores actualizados $w(1 \quad 0.625), b= -0.125$.

\subsubsection*{Muestra 2}
Tenemos $x_1 = 0$, $x_2 = 1$ y $y=0$, así que

\begin{align*}
    y_{pred} &= 1\cdot 0 + 0.625 \cdot 1 - 0.125 = 0.5 \\
    error &= 0 - 0.5 = -0.5 \\
    w_1 &= 1 + 0.5 \cdot (-0.5) \cdot 0 = 1 \\
    w_2 &= 0.625 + 0.5 \cdot (-0.5) \cdot 1 = 0.375 \\
    b &= -0.125 + 0.5 \cdot (-0.5) = -0.375
\end{align*}
Se tienen los valores actualizados $w(1 \quad 0.375), b= -0.375$.

\subsubsection*{Muestra 3}
Tenemos $x_1 = 1$, $x_2 = 0$ y $y=0$, tenemos
\begin{align*}
    y_{pred} &= 1 \cdot 1 + 0.375 \cdot 0 - 0.375 = 0.625 \\
    error &= 0 - 0.625 = -0.625 \\
    w_1 &= 1 + 0.5 \cdot (-0.625) \cdot 1 = 0.6875 \\
    w_2 &= 0.375 + 0.5 \cdot (-0.625) \cdot 0 = 0.375 \\
    b &= -0.375 + 0.5 \cdot (-0.625) = -0.6875
\end{align*}
Se tienen los valores actualizados $w(0.6875 \quad 0.375), b= -0.6875$.

\subsubsection*{Muestra 4}
Tenemos $x_1 = 1$, $x_2 = 1$ y $y=1$, tenemos
\begin{align*}
    y_{pred} &= 0.6875 \cdot 1 + 0.375 \cdot 1 - 0.6875 = 0.375 \\
    error &= 1 - 0.375 = 0.625 \\
    w_1 &= 0.6875 + 0.5 \cdot (0.625) \cdot 1 = 1 \\
    w_2 &= 0.375 + 0.5 \cdot (0.625) \cdot 1 = 0.6875 \\
    b &= -0.6875 + 0.5 \cdot (0.625) = -0.375
\end{align*}
La iteración 2 finaliza con los valores $w(1 \quad 0.6875), b= -0.375$.

\subsection*{Iteración 3}
\subsubsection*{Muestra 1}
Tenemos $x_1 = 0$, $x_2 = 0$ y $y=0$, tenemos
\begin{align*}
    y_{pred} &= 1\cdot 0 + 0.6875 \cdot 0 - 0.375 = -0.375 \\
    error &= 0 - (-0.375) = 0.375 \\
    w_1 &= 1 + 0.5 \cdot (0.375) \cdot 0 = 1 \\
    w_2 &= 0.6875 + 0.5 \cdot (0.375) \cdot 0 = 0.6875 \\
    b &= -0.375 + 0.5 \cdot (0.375) = -0.1875
\end{align*}
Se tienen los valores actualizados $w(1 \quad 0.6875), b= -0.1875$.

\subsubsection*{Muestra 2}
Tenemos $x_1 = 0$, $x_2 = 1$ y $y=0$, así que
\begin{align*}
    y_{pred} &= 1\cdot 0 + 0.6875 \cdot 1 - 0.1875 = 0.5 \\
    error &= 0 - 0.5 = -0.5 \\
    w_1 &= 1 + 0.5 \cdot (-0.5) \cdot 0 = 1 \\
    w_2 &= 0.6875 + 0.5 \cdot (-0.5) \cdot 1 = 0.4375 \\
    b &= -0.1875 + 0.5 \cdot (-0.5) = -0.4375
\end{align*}
Se tienen los valores actualizados $w(1 \quad 0.4375), b = -0.4375$.

\subsubsection*{Muestra 3}
Tenemos $x_1 = 1$, $x_2 = 0$ y $y=0$, tenemos
\begin{align*}
    y_{pred} &= 1 \cdot 1 + 0.4375 \cdot 0 - 0.4375 = 0.5625 \\
    error &= 0 - 0.5625 = -0.5625 \\
    w_1 &= 1 + 0.5 \cdot (-0.5625) \cdot 1 = 0.71875 \\
    w_2 &= 0.4375 + 0.5 \cdot (-0.5625) \cdot 0 = 0.4375 \\
    b &= -0.4375 + 0.5 \cdot (-0.5625) = -0.71875
\end{align*}
Se tienen los valores actualizados $w(0.71875 \quad 0.4375), b= -0.71875$.

\subsubsection*{Muestra 4}
Tenemos $x_1 = 1$, $x_2 = 1$ y $y=1$, tenemos
\begin{align*}
    y_{pred} &= 0.71875 \cdot 1 + 0.4375 \cdot 1 - 0.71875 = 0.4375 \\
    error &= 1 - 0.4375 = 0.5625 \\
    w_1 &= 0.71875 + 0.5 \cdot (0.5625) \cdot 1 = 1 \\
    w_2 &= 0.4375 + 0.5 \cdot (0.5625) \cdot 1 = 0.71875 \\
    b &= -0.71875 + 0.5 \cdot (0.5625) = -0.4375
\end{align*}
La iteración 3 finaliza con los valores $w(1 \quad 0.71875), b= -0.4375$.

\subsection*{Verificación}
Finalmente, verificaremos que el modelo clasifica a todas las muestras partiendo de:
\[y_{pred} = 1 \cdot x_1 + 0.71875 \cdot x_2 - 0.4375\]
Entonces
\begin{center}
\begin{tabular}{cccc}
    \toprule
    $x_1$ & $x_2$ & $y$ & $y_{pred}$ \\
    \midrule
    0   & 0   & 0   & -0.4375 \\
    0   & 1   & 0   & 0.28125 \\
    1   & 0   & 0   & 0.5625 \\
    1   & 1   & 1   & 1.28125 \\
     \bottomrule
    \end{tabular}
\end{center}

Si tomamos un umbral de $0.5$, se tiene que el modelo clasifica correctamente a $3$ de las $4$ muestras.